%----------------------------------------------------------------------------------------
% Tailored CV — Wissenschaftliche*r Mitarbeiter*in, Tumor Immunology / Gastric Cancer
% Klinik VTG, Universitätsklinikum Carl Gustav Carus Dresden — Ref. 1776
%----------------------------------------------------------------------------------------

\documentclass[10pt,a4paper,sans]{moderncv}

\usepackage[utf8]{inputenc}
\usepackage[T1]{fontenc}

\moderncvstyle{classic}
\moderncvcolor{blue}

\usepackage[scale=0.78]{geometry}
\usepackage{ragged2e}

%----------------------------------------------------------------------------------------

\firstname{Kundai Farai}
\familyname{Sachikonye}

\address{Auf der Lichtung 19}{82178 Puchheim, Germany}
\mobile{(+49) 1626386344}
\email{kundai.sachikonye@bitspark.com}
\homepage{github.com/fullscreen-triangle}{github.com/fullscreen-triangle}

%----------------------------------------------------------------------------------------

\begin{document}

\makecvtitle

%----------------------------------------------------------------------------------------
%	EDUCATION
%----------------------------------------------------------------------------------------

\section{Education}

\cventry{2019--2021}{PhD candidate, Computational Lipidomics}{Technische Universität München / Bayerisches Forschungsinstitut}{}{}{
  Multi-omics statistical modelling, mass spectrometry-based molecular profiling,
  federated learning for distributed clinical analysis. Departed to pursue independent research.
}
\cventry{2018--2019}{MSc Computational Biology}{Universität Tübingen}{}{}{
  Computational immunology, sequence analysis, multi-omics integration,
  statistical bioinformatics.
}
\cventry{2014--2017}{MSc Pharmaceutical Biotechnology}{HAW Hamburg}{}{}{
  Cell biology and culture techniques, protein biochemistry, immunological
  assay development, bioanalytical methods.
}
\cventry{2011--2014}{BSc Biochemistry and Cell Biology}{Jacobs University Bremen}{}{}{
  Molecular biology, cell biology, analytical biochemistry, cancer biology.
}

%----------------------------------------------------------------------------------------
%	RELEVANT MANUSCRIPTS
%----------------------------------------------------------------------------------------

\section{Relevant Manuscripts}

\cvitem{2025}{
  \textbf{On the Thermodynamic Consequences of Categorical Completion on Biological Systems.}
  MHC binding affinity correlates with categorical richness ($r=0.73$, $p<10^{-12}$,
  $n=2{,}847$ IEDB peptides). AUC 0.81 vs.\ NetMHCpan 0.68.
  \textit{Manuscript.}
}

\cvitem{2025}{
  \textbf{On Disease Epistemology in Blind Receiver.}
  Loop-holonomy self-consistency diagnostic; AUC$=1.00$; 25/25 validation experiments;
  sparse $\ell_1$ LP for personalised therapeutic target selection.
  \textit{Manuscript. AIMe Registry / TUM Weihenstephan.}
}

%----------------------------------------------------------------------------------------
%	COMPUTATIONAL TOOLS FOR TUMOR IMMUNOLOGY
%----------------------------------------------------------------------------------------

\section{Computational Tools for Tumor Immunology}

\cvitem{syndrome}{
  \textbf{Neoantigen Prediction, MHC Binding, and Immune Escape Analysis.}
  Spectral DFT on three physicochemical channels (hydropathy, volume, charge);
  36-dimensional cosine similarity $\rho$; $\mathcal{O}(cL\log L)$.
  \textbf{Neoantigen:} dual threshold $\rho_\mathrm{MHC}>0.72$ (MHC presentation)
  + $\Delta\rho_\mathrm{self}>0.15$ (self-distinction); ranks neoantigens by
  immunogenic potential in patient tumor mutational landscape.
  \textbf{Escape:} $E(v)=-\max\rho(v,\mathrm{Ab})+0.80\cdot\rho(v,\mathrm{rec})$;
  systematic screen of all single-aa substitutions for variants that escape
  T cell / antibody recognition while retaining receptor function.
  \textbf{MHC binding:} $\rho\geq0.72$ threshold; patient-specific HLA allele scanning.
  Applicable to gastric cancer neoantigen identification and tumor immune escape modeling.
  Live: \href{https://syndrome-seven.vercel.app/tools}{syndrome-seven.vercel.app/tools}.
  \href{https://github.com/fullscreen-triangle/syndrome}{github.com/fullscreen-triangle/syndrome}
}

\cvitem{syndrome disease model}{
  \textbf{T Cell State and Tumor Microenvironment Trajectory.}
  Disease / cell-state vector $\mathbf{D}\in[0,1]^8$ across eight oscillator
  dimensions (ATP synthesis coherence, genetic expression state, immune activation,
  circadian regulation, etc.); $\mathcal{O}(N\log N)$ trajectory to target state;
  sparse $\ell_1$ LP for minimum-intervention target selection.
  Applicable to modelling T cell exhaustion trajectories and identifying
  metabolic rescue or checkpoint blockade combination targets.
  Loop-holonomy self-consistency: AUC$=1.00$, 25/25 experiments.
}

\cvitem{lavoisier}{
  \textbf{Spectral Proteomics and Metabolomics Analysis.}
  Dual-pipeline MS analysis: numerical + CNN (PyTorch); 98.9\,\% NIST accuracy;
  Ray distributed; $>$100\,GB dataset support.
  Applicable to metabolomic profiling of T cell metabolic phenotype (OxPhos vs.\
  glycolytic flux) from tumor-infiltrating lymphocyte samples.
  \href{https://github.com/fullscreen-triangle/lavoisier}{github.com/fullscreen-triangle/lavoisier}
}

\cvitem{gospel}{
  \textbf{Genomic Variant Analysis and Multi-Omics Integration.}
  $10^6$ variants/second VCF processing; RNA-seq integration; Bayesian networks,
  GMMs. Validated: 1000~Genomes, gnomAD~v3.1.2, UK~Biobank.
  Applicable to tumor mutational burden analysis and somatic mutation landscape
  for neoantigen candidate identification in gastric cancer.
  \href{https://github.com/fullscreen-triangle/gospel}{github.com/fullscreen-triangle/gospel}
}

%----------------------------------------------------------------------------------------
%	EXPERIENCE
%----------------------------------------------------------------------------------------

\section{Experience}

\cventry{2021--present}{Independent AI \& Computational Biology Developer}{Bitspark GmbH}{}{}{
  Computational immunology tools (neoantigen, MHC binding, immune escape),
  disease state trajectory modelling, multi-omics analysis pipelines,
  proteomics and metabolomics data analysis. All systems publicly deployed.
}

\cventry{2019--2021}{Computational Lipidomics}{TUM / Bayerisches Forschungsinstitut}{Munich}{}{
  MS-based lipidomics: molecular profiling, peak detection, ion annotation
  (LIPID MAPS, LipidBlast), multi-omics integration, federated clinical data.
  Python, Java.
}

\cventry{2017--2018}{Glycomics and Proteomics}{Universitätsklinikum Hamburg-Eppendorf}{}{}{
  Automated LC-MS/MS and MALDI-TOF pipeline; protein sample preparation;
  SOP development and reproducibility validation.
}

\cventry{2013}{Cancer Biomarker Discovery}{University Medical Center Groningen}{}{}{
  UPLC-MS/MS shotgun proteomics for cancer biomarker discovery;
  sample preparation protocol development; statistical validation.
}

%----------------------------------------------------------------------------------------
%	TECHNICAL SKILLS
%----------------------------------------------------------------------------------------

\section{Technical Skills}

\cvitem{Tumor Immunology (Computational)}{
  Neoantigen prediction; MHC/HLA binding; TCR--pMHC spectral complementarity;
  immune escape variant screening; T cell exhaustion modelling;
  tumor mutational burden analysis; somatic variant calling (VCF, PLINK);
  flow cytometry data analysis (FlowJo-compatible, Python-based)
}

\cvitem{Cell Biology (Training)}{
  Cell culture; protein extraction and Western blot; ELISA; immunostaining;
  histological staining; basic flow cytometry acquisition
}

\cvitem{Analytical / MS}{
  LC-MS/MS, MALDI-TOF, UPLC-MS/MS; proteomics and lipidomics pipelines;
  LIPID MAPS, HMDB, NIST; spectral data analysis (numerical + CNN)
}

\cvitem{ML / Bioinformatics}{
  Python (PyTorch, scikit-learn, Ray); R; Rust;
  Bayesian networks; GMMs; sparse $\ell_1$ regularisation;
  multi-omics integration; RNA-seq quantification
}

%----------------------------------------------------------------------------------------
%	LANGUAGES
%----------------------------------------------------------------------------------------

\section{Languages}

\cvitemwithcomment{English}{Native / Fluent}{}
\cvitemwithcomment{German}{Conversational}{Living in Germany since 2011}

%----------------------------------------------------------------------------------------

\renewcommand{\listitemsymbol}{-~}

\end{document}
